\chapter{Numerical utilities}

\name{} library includes various numerical code snippets, gathered
into the sublibrary \texttt{ntools}.
The routines concerns various aspects of numerical programming, that
can be useful in various situations.
By convention, the routines names always start with the prefix \verb|ntools_|.
We also try to benefit form preexisting conventions in packages as
\blas{} and \verb|LAPACK|\index{LAPACK@LAPACK} in order to make naming
schemes as consistent as possible.
The presented conventions will also be used in the subsequent chapters.

\section{Constants and special values}

Mathematical constants are represented in the file \verb|constants.h|,
and are summarized in Table~\ref{tab:constants}.
\begin{table}[htb]
\begin{small}
\begin{tabularx}{\textwidth}{lll}
Symbol & Value & Numerical approximation \\
\hline
\verb|M_E| & $e$ & 2.71828182845904523536028747135 \\
\verb|M_EULER| & $\Gamma'(1)$ (Euler constant\index{Euler constant})&
0.5772156649015328606065 \\
\verb|M_LN2| & $\ln(2)$ & 0.69314718055994530941723212146 \\
\verb|M_LNPI| & $\ln(\pi)$ & 1.14472988584940017414342735135 \\
\verb|M_LN10| & $\ln(10)$ & 2.30258509299404568401799145468 \\
\verb|M_LOG2E| & $\log_2(e)$ & 1.44269504088896340735992468100 \\
\verb|M_LOG10E| & $\log_{10} (e)$ & 0.43429448190325182765112891892 \\
\verb|M_PI| & $\pi$ & 3.14159265358979323846264338328 \\
\verb|M_PI_2| & $\pi/2$ & 1.57079632679489661923132169164 \\
\verb|M_PI_4| & $\pi/4$ & 0.78539816339744830961566084582 \\
\verb|M_1_PI| & $1/\pi$ & 0.31830988618379067153776752675 \\
\verb|M_2_PI| & $2/\pi$ & 0.63661977236758134307553505349 \\
\verb|M_SQRT1_2| & $\sqrt{1/2}$ & 0.70710678118654752440084436210 \\
\verb|M_SQRT2| & $\sqrt{2}$ & 1.4142135623730950488 \\
\verb|M_SQRT3| & $\sqrt{3}$ & 1.73205080756887729352744634151 \\
\verb|M_SQRTPI| & $\sqrt{\pi}$ & 1.77245385090551602729816748334 \\
\verb|M_TWO_SQRTPI| & $2/\sqrt{\pi}$ & 1.12837916709551257389615890312 \\
\verb|M_TWO17| & $2^{17}$ & 131072.0 \\
\verb|M_TWO53| & $2^{53}$ & 9007199254740992.0 \\
\hline
\end{tabularx}
\label{tab:constants}
\caption{Mathematical constants}
\end{small}
\end{table}

\begin{verbatim}
double nt_infinity()
double nt_nan()
\end{verbatim}

\section{Geometric routines}

\subsection{Cartesian and polar coordinates}

Given a two-dimensional vector $a$, the two classical coordinates system are the cartesian and polar ones.
In the Cartesian system, $a$ is expressed as a linear combination of the canonical vectors $e_1 = (1,0)$ and $e_2=(0,2)$ (defining the $x$- and $y$-axis):
\[
a = x e_1 + y e_2.
\]
Therefore, we can write $a$ as the couplet $(x, y)$.
In the polar system, $(r, \theta)$, where $r$ is the radial distance from the origin, and $\theta$ is the counterclockwise angle from the $x$-axis.
Storing the Cartesian coordinates in a vector, we have the two following functions, that perform conversion from Cartesian to polar coordinates and from polar to Cartesian coordinates, respectively:
\begin{verbatim}
  void nt_cartesian2polar(double *x, double r, double theta);
  void nt_polar2cartesian(double r, double theta, double *x);
\end{verbatim}

\section{Derivatives}
\index{derivatives}

\subsection{Finite approximations}

A common way to compute approximate first and second derivatives is to use finite-difference
approximations.\index{finite difference}

\subsubsection{Finite-difference gradients}

\section{Linear algebra}

\subsection{Conventions}
\label{sub:laconventions}

Each real vector $\boldsymbol{x}$ is represented as a double array
whose consecutive elements are separated using a fixed increment, as
illustrated in Figure~\ref{fig:vector}.
\begin{figure}
\[
\boldsymbol{x}: \underbrace{x_0 \sqcup \sqcup \ldots
  \sqcup}_{\mbox{\texttt{incx}}}x_1 \ldots
\]
\label{fig:vector}
\caption{Vector representation}
\end{figure}

When manipulating matrices,
we follow the \lapack{} convention (see for instance
Hogben~\cite{Hogb06}, Chapter~75), and denote the number of rows by $m$
and the number of columns by $n$.
We also define the leading dimension\index{dimension!leading} as the
first dimension of a two-dimensional array (as opposed to the
dimension of the matrix stored in that array).
Since arrays are stored by rows in C, the leading dimension is often the
number of columns, but it can greater than this number, allowing
convenient abstraction to access submatrices.
The basic memory allocation and deallocation are
\begin{verbatim}
double **nt_matrix_new(int m, int n)
double **nt_matrix_realloc(double **A, int m, int n)
void nt_matrix_free(double **A)
\end{verbatim}

Various additional functions offer basic operations on matrices.
The function
\begin{verbatim}
void nt_matrix_zeros(int m, int n, double *A, int ldA)
\end{verbatim}
set all elements of the $m \times n$ leading principal submatrix of A to 0.0.
\verb|ldA| is the leading dimension of $A$.
This submatrix can also be set to the $m \times n$ identity matrix
\index{identity!matrix} with the function
\begin{verbatim}
void nt_matrix_identity(int m, int n, double *A, int ldA)
\end{verbatim}
These functions implicitely assume that the matrix is row-oriented, as
said before, but this is not the case with every
language. In particular, Fortran stores matrices by columns.
Since many numerical routines are available in Fortran, and can be
easily linked to C code, it can be useful to specify how matrices are
stored.
This can can be done using specialized functions, whose names start by
\verb|nt_ma| instead of \verb|nt_matrix|, and whose first argument is
the matrix orientation identifier.
It has the specific type \verb|enum CBLAS_ORDER|, and can take one of
the following constant values:
\begin{enumerate}
\item
\verb|NT_MATRIX_ROW| or \verb|CblasRowMajor| for a row-oriented matrix;
\item
\verb|NT_MATRIX_COLUMN| or \verb|CblasColMajor| for a column-oriented matrix.
\end{enumerate}
The constants
\verb|CblasRowMajor|\index{CblasRowMajor@\texttt{CblasRowMajor}} and
\verb|CblasColMajor|\index{CblasColMajor@\texttt{CblasColMajor}} are
proposed in order to be consistent with \blas{} (see
Section~\ref{sec:requirements}) conventions.
The corresponding functions list is
\begin{verbatim}
double **nt_ma_new(const enum CBLAS_ORDER Order, int m, int n)
void nt_ma_diag(const enum CBLAS_ORDER Order,
		const int m, const int n, double *A, const int ldA,
		double alpha)
void nt_ma_identity(const enum CBLAS_ORDER Order,
		    const int m, const int n, double *A, const int ldA)
void nt_ma_scale(const enum CBLAS_ORDER Order, const int m, const int n,
		 double alpha, double *A, int ldA)
void nt_ma_zeros(const enum CBLAS_ORDER Order,
		 const int m, const int n, double *A, const int ldA)
\end{verbatim}

\subsection{I/O routines}

These routines allow simple input/output management for some
mathematical constructions.
\begin{verbatim}
void nt_vect_print(FILE *out, char *name,
                   double *v, int n)
\end{verbatim}
\begin{verbatim}
void nt_matrix_print(FILE *out, char *name,
		     double **A, int m, int n)
\end{verbatim}

\subsection{Modular arithmetic}
\index{arithmetic!modular}

Let $m>0$ be a natural number, and consider the set ${\zit}_m
\overset{def}{=}\{ 0,1,\ldots,m-1\}$.
For $a,b\in {\mathbb{Z}}_n$ define $a+b$ and $ab$ by performing these
operations in $\mathbb{Z}$, and the subtracting multiples of $m$ until
the result is in ${\mathbb{Z}}_m$.
In other words, form the sum and product of $a$ and $b$ as integers, divide
them by $m$ and take the remainders.
These operations are called the addition and multiplication modulo
$m$.
Under these relations, we have that two integers $a$ and $b$ are said
to be congruent modulo $m$\index{congruence} if their difference is a
multiple of $m$ (in other terms the remainder of the integer division
$(a-b)/m$ is equal to 0).
We will write $a = b \mod m$.

\subsubsection{Multiplication: $ax \mod m$}

Mupltiplication modulo $m$ is an important tool in various
application, for instance random numbers generations.
The main difficulty is that with large $m$, a crude implementation of
$ax \mod m$ could lead to overflows.
We examine two ways to circumvent this problem.

\paragraph{Approximate factorization}
\index{factorization!approximate}

This method is also know as \Index{Schrage}'s method, and is valid if
$a^2 < m$ or $a = \lfloor m/i \rfloor$, where $i^2 < m$.
We work with integer only, and precompute ${q} = \lfloor m/a \rfloor$
and ${r} = m\mod a$.
We then 
\[
  y \leftarrow \lfloor x / q \rfloor; \quad
  x \leftarrow a (x - yq) - yr; \quad
  \mbox{If } x < 0 \mbox{ then } x \leftarrow x + m.
\]%
We indeed have that
\begin{eqnarray*}
  a x \mod m &=&  (ax - \lfloor x/q\rfloor m) \mod m \\
      &=&  (ax - \lfloor x/q\rfloor (a q + r)) \mod m \\
      &=&  (a (x - \lfloor x/q\rfloor q) - \lfloor x/q\rfloor r) \mod m \\
      &=&  (a (x \mod q) - \lfloor x/q\rfloor r) \mod m.
\end{eqnarray*}
Under our assumptions, it is easy to show that all intermediate
quantities stay between $-m$ amd $m$.

The following function computes $ax mod m$ and returns the result of
the operation
\begin{verbatim}
long nt_modular_schrage(long a, long x, long m)
\end{verbatim}
While we propose a generic function, it is usally more efficient to
directly implement this technique inline, by precomputing $q$ and $r$,
as explined.

\paragraph{Floating-points computation}

If we assume that we work with double precision numbers coded on 64
bits, with a mantissa represented by 53 bits, the presented technique
is valid if $a.m < 2^{53}$, since the we can represent all integers
betweem 0 and $2^{53}-1$ with no rounding errors.
The following algorithm produces the desired result.
\begin {eqnarray*}
   && \mbox{double } m, a, x, v; \quad  \mbox{int } k; \\
   && v = a.x; \quad
    k = \lfloor v / m \rfloor; \quad
    x = v - k.m;
\end {eqnarray*}
The following function computes $ax + y mod m$:
\begin{verbatim}
double nt_modular_daxpy(double a, double x, double c, double m)
\end{verbatim}
The name \verb|daxpy| is used as reference to the \blas{} operation $ax+y$.

\subsubsection{Power}

The multiplication operation is much simpler when $x$ is equal or more
generally to a power of $a$.
The following routine performs the operation $a^e \mod{} m$.
\begin{verbatim}
long nt_modular_power(long a, long e, long m);
\end{verbatim}
\section{Primes numbers\index{prime}}

More advanced and efficient algorithm exist, but their implementation
goes far over the purpose of this library.
We only propose simple routines, that you should be sufficient for
numerous applications, when only a few prime numbers are required.
If you require a general but fast prime numbers search, we recommand
the use of Atkin's sieve~\cite{AtkiBern04}.

\subsection{Successive divisions}

The simple methods consists to enumerate naturals and check if their
primes by testing if they can be divided by a previously computed
prime.
The method is less memory consuming that the Eratosthene sieve
presented below, but is requires more CPU time.
You can compute the first $n$ prime number, with the following function:
\begin{verbatim}
void nt_primes_simple(unsigned int n, unsigned int *primes)
\end{verbatim}
The array \verb|primes| has to be memory allocated.
On output, \verb|primes[i]| will containes the $(i-1)^{\mbox{th}}$
prime.
As an illustration, the 10 first prime numbers are 2, 3, 5, 7, 11, 13,
17, 19, 23, 29.

\subsection{Eratosthene sieve\index{Eratosthene!sieve}}

\begin{verbatim}
void nt_primes_eratosthene(unsigned int n, unsigned int *primes)
\end{verbatim}

\section{Special functions}

Various functions are directly supported in \texttt{ORATIO} as we can
frequently encounter them, and either implementaion is not trivial,
either we can try to optimize their efficiency.

\subsection{B-Splines\index{spline}}

Various utility functions are defined to deal with B-splines.
A B-spline function of degree $p$ is a polynomial function of degree
$p$, defined on the interval $[a, b]$, that can be expressed as a linear
combination of $n+1$ basis functions $N_{i,p}(u)$, as follows:
\begin{equation}
C(u) = \sum_{i=0}^n P_i N_{i,p}(u) .
\label{eq:spline_expression}
\end{equation}
The coefficients $P_0, P_1, \ldots, P_n$ are called the control
points, and $u$ is the knot vector $(u_0 = a, u_1, \ldots, u_m
=b)$.
The basis functions can be constructed by recurrence on the degree $p$:
\[
N_{i, 0} =
\begin{cases}
1 & \mbox{if } u \in [u_i, u_{i+1}), \\
0 & \mbox{otherwise.}
\end{cases}
\]
and
\[
N_{i,p} = \frac{u-u_i}{u_{i+p}-u_i} N_{i, p-1}(u) +
\frac{u_{i+p+1}-u}{u_{i+p+1}-u_{i+1}} N_{i+1, p-1}(u),
\]
so that $n$ is equal to $m-p-1$.

There are several types of knot vectors, but one especially convenient
is the nonperiodic (or clamped or open) knot vector, which has the form 
\[
U = \lbrace \underbrace{a,\ldots,a}_{p+1}, u_{p+1}, \ldots, u_{m-p-1},
\underbrace{b,\ldots,b}_{p+1} \rbrace,
\]
that is the first and last knots have multiplicity $p+1$.
It is possible to show that the function $C(u)$ is $p-1$ continuously
differentiable.

The implementation proposed in the library is based on Piegl and
Tiller~\cite{PiegTill96}.
A space of B-splines is represented with the type \verb|NT_BSpline|, defining the set
of a basis function and the knot vector.
A new B-spline space can be created using the function
\begin{verbatim}
NT_BSpline *nt_bsplines_new(int p, int nk, double *knots, int inck)
\end{verbatim}
where, as before, \verb|p| is the degree of the spline, \verb|nk| is
the dimension of the knots vector pointed by \verb|knots|, with the
increment \verb|inck|, which should be set to 1 is the elements or
consecutive.
The allocated memory can further be free using
\begin{verbatim}
void nt_bsplines_free(NT_BSpline *bspline)
\end{verbatim}
This function internally creates the knots vector $U$ (with multiplicity $p$+1 for the extreme points)
using the following function, that returns the number of components in U, i.e. \verb|nk| plus 2$p$.
\begin{verbatim}
int nt_bsplines_knots_create(int p, int nk, double *knots,
                             int *nU, double **U)
\end{verbatim}
The $U$ dimension can also be recovered with the function
\begin{verbatim}
int nt_bsplines_knots_dimension(int p, int nk)
\end{verbatim}
and a vector of appropriate size can be created with the function
\begin{verbatim}
double *nt_bsplines_knots_allocate(int p, int nk)
\end{verbatim}
A knot vector with all multiplicities equal to 1, and $n$ components, can be adjusted as an open knot vector with the function below.
\begin{verbatim}
int nt_bsplines_knots_initialize(int p, int *n, double **U,
                                 int realloc)
\end{verbatim}
The variable \verb|realloc| has to be set to 1 if a memory reallocation has to be performed, and 0 otherwise. On output, $U$ is the open knot vector, and $n$ is the number of its components, also returned by the function.

The dimension $n$ in \eqref{eq:spline_expression}, 
given the degree $p$ of the spline basis functions, and the dimension $m+1$ of the knot vector
with multiplicities set to 1, can be retrieved with the function
\begin{verbatim}
int nt_bsplines_space_dimension(int p, int m)
\end{verbatim}

In order to define a particular B-Spline in such a space, the user has to specify the vector of control
points $P$.
The B-Spline can then be evaluated at a point $u$ with the function
\begin{verbatim}
double nt_bsplines_evaluate(NT_BSpline *b, double *P,
                            double u)
\end{verbatim}
This function is the short form of the routine below, where \verb|n| is the dimension of the open knot vector \verb|U|, \verb|p| is the spline degree, and \verb|N| is a vector of size at least equal to \verb|p|+1, that will be used to recursively evaluate the basis functions:
\begin{verbatim}
double nt_bsplines_curve_point(int n, int p, double u, double *U,
                               double *P, double *N)
\end{verbatim}

%\begin{verbatim}
%int nt_bsplines_find_span(int n, int p, double u, double *U)
%\end{verbatim}

%\begin{verbatim}
%void nt_bsplines_basis_func(int i, double u, int p,
%                            double *U, double *N)
%\end{verbatim}

%\begin{verbatim}
%double nt_bsplines_one_basis_func(int p, int m, double *U,
%                                  int i, double u)
%\end{verbatim}

\subsection{Power}

\subsubsection{Integer power}

The standard function \verb|pow|, in the mathematical library, only consider $x^y$, where $x$ and $y$ are real (with the restriction that $y$ must be an integer if $x$ is negative.
When $y$ is some integer, it is possible to simplify the computation
by only performing multiplications, whose number can be limited.
The function is similar to \verb|pow| as it takes two arguments, $x$ and $y$, wehere $y$ is restricted to be an integer, and return the value $x^y$:
\begin{verbatim}
double nt_ipow(double x, int y)
\end{verbatim}
On a Macbook Pro 2.0Ghz, and compiler, gcc-4, our routine appeared
significantly faster (by a factor of 2 to 3) than the standard
function for simple tests.

\subsection{Factorial}
\index{factorial}

The value of $n!$, where $n$ is a nonnegative integer, can be obtained with the function
\begin{verbatim}
double nt_factorial(int n)
\end{verbatim}

\subsubsection{Powers of 2}
\index{power}

It is possible to precalculate an array containing powers of 2 using
to function
\begin{verbatim}
void nt_two_powers_array(int n, unsigned int *p)
\end{verbatim}
Then \verb|p[i]| will be equal to $2^i$, as long a $i$ is not too
large compared to the integer machine precision.

It is also possible to compute the power of 2 closest but inferior to
some integer $x$, i.e $i$, such that $2^i <= x$, and for all $j > i$, $2^j > x$.
\begin{verbatim}
int nt_two_power_floor(int x)
\end{verbatim}
Similarly, the function
\begin{verbatim}
int nt_two_power_ceil(int x)
\end{verbatim}
returns $i$, where $2^i >= x$, and for all $j < i$, $2^j < x$.

\subsection{Radical inverse}
\index{radical inverse}

Let $i \in \nit$, and $b \in \nit_0$.
The $b$-ary expansion of $i$ is
\[
i = a_0 + a_1b + \ldots + a_{k-1}b^{k-1},
\]
and its radical inverse in base $b$ is
\[
\psi_b(i) = a_0b^{-1} + a_1b^{-2} + \ldots + a_{k-1}b^{-k}.
\]
The value of $\psi_b(i)$ can be computed using the function
\begin{verbatim}
double nt_radical_inverse(int b, int i);
\end{verbatim}
